% Lösungen Einheit 3 von 4 – Parameter und Transformationen (zusammenbau v0.1)
\einheitenkopf{Einheit 3 von 4 · Parameter und Transformationen}

%% TODO Lösung der Erklärzeile 18g) fehlt in der Bank
\erg{18}{a) die Höhe (a) \quad b) 4 \quad c) 2,5 \quad d) 6 \quad e) 1,5 \quad f) 8 \quad g) \ldots \quad h) Grenzen −5 und 5: alle Werte von −5 bis 5 \quad i) a = 2,5 (halber Abstand zwischen Hochpunkt 2,5 und Tiefpunkt −2,5) \quad j) an der waagerechten Achse gespiegelt: die Welle geht zuerst nach unten, T(90|−2) \quad k) Periode = 360° : 3 = 120° \quad l) Periode 90°, b = 360° : 90° = 4 \quad m) Welle zwischen −2 und 2 mit Periode 120°: H(30|2), T(90|−2), Nullstellen alle 60° \quad n) a = 4; Periode 240°; b = 360° : 240° = 1,5; also y = 4 · sin(1,5 · x) \quad o) y = sin(4 · x) \quad p) höher: sie schwingt zwischen −4 und 4, die Breite bleibt gleich \quad q) um 2 nach oben; Kleinstwert = 1, Größtwert = 3 \quad r) um 90° nach links: y = sin(x + 90°) \quad s) Periode 90°; Nullstelle bei x = 45° (halbe Periode), Hochstelle bei x = 22,5° (Viertelperiode) \quad t) a) y = 3,5 · sin(3 · x) (Amplitude 3,5, Periode 120°); b) Welle zwischen −2 und 2 mit Periode 240°: H(60|2), T(180|−2)}

18m)

\begin{ksys}[xmin=0,xmax=360,ymin=-3,ymax=3,xstep=30,ystep=1,klein]\funktion{2*sin(3*\x)}{f}\end{ksys}


18t)

\begin{ksys}[xmin=0,xmax=480,ymin=-3,ymax=3,xstep=60,ystep=1,klein]\funktion{2*sin(1.5*\x)}{f}\end{ksys}

\erg{19}{a) Amplitude = 2 bei beiden; Periode von g = 120°, von f = 360° – g ist dreimal so schnell}
\erg{20}{a) Fehler: mit b malgenommen statt geteilt. Richtig: Periode = 360° : 4 = 90°, die Welle wird schmaler \quad b) Bei größerem b läuft der Punkt schneller um den Einheitskreis: der Vollkreis ist schon nach einem kleineren Winkel x durchlaufen, die Periode wird kürzer.}
